\chapter{Einleitung}
    
Physik ist die Grundlage für das Verständnis technischer Systeme im Maschinenbau und in der Systemtechnik.

Dieses Skript dient als zentrales Medium für den Physikunterricht im 3. Semester für Maschinentechniker Stufe HF an der Inovatech. 
Aufbauend auf den Grundkenntnissen der Mathematik und dem Physikskript von Semester 2 \cite{schnidrig_physik_2026}
vermittelt es die wichtigsten physikalischen Prinzipien, 
die für die Entwicklung und Optimierung moderner Technologien essenziell sind. Von Statik über Dynamik bis hin zur 
Strömungslehre werden die Inhalte möglichst anschaulich und praxisnah aufbereitet.

Wo auch immer möglich werden Beispiele und Aufgaben an reellen physikalischen Systemen gezeigt.
Diese sind meist stark vereinfacht und viele Einflüsse werden vernachlässigt, allerdings ist dies die 
Kunst für alle Stufen der Physik. Die Realität ist beliebig komplex, 
wer sie verständlich und isoliert vereinfachen kann, hat einen grossen Vorteil.

Nutzen Sie dieses Skript als Leitfaden, 
um den Einstieg in die Physik zu meistern und Ihre technischen Fähigkeiten zu vertiefen. 
Die behandelten Themen bieten eine solide Basis, die Sie je nach Interesse und Bedarf 
selbstständig vertiefen können. Literatur ist für alle Themengebiete im Überfluss erhältlich. 
Ziel des Unterrichts ist es, Ihnen die wichtigsten Werkzeuge 
an die Hand zu geben, um sich in neue Themengebiete einzuarbeiten und diese erfolgreich anzuwenden.

\begin{figure}[H]    
    \centering          %     left botm right top
    \includegraphics[clip, trim=0cm 0cm 0cm 0cm,width=.3\textwidth]{physics_logo.jpg}
    \label{fig:einleitungPhysicsLogo}
\end{figure}

\newpage

\section{Organisation Unterricht}

Der Unterricht wird über OpenOlat organisiert. Der Unterricht ist als einfache Version des 
Flipped-Classrooms aufgebaut. Auf OpenOlat wird pro Unterrichtsblock ein Themenbereich oder Teile eines
Themenbereichs zugewiesen. Dieser ist vor dem jeweiligen Unterrichtsblock vorzubereiten.

In der ersten Lektion des Unterrichtblocks wird in einem kurzen Wrap-up der vorbereitete Stoff wiederholt.
Anschliessend werden allfällige Fragen beantwortet, der Rest des Unterrichts wird für das Lösen von Aufgaben genutzt.

Das Skript kann sich während des Semesters ändern, bitte achten Sie deshalb darauf, dass Sie auf
OpenOlat jeweils die aktuellste Version benutzen.

\section{Zu dem Skript}

Das Skript ist in die Themenbereiche aufgeteilt:

\begin{itemize}
    \item Statik 2
    \item Kinematik 2
    \item Hydrostatik
    \item Kontinuität und Bernoulli
    \item Schwingungen und Wellen 2
\end{itemize}

Zu jedem Themenbereich sind Übungen und Lösungen in den Anhängen \ref{app:A_Übungen} und \ref{app:B_Lösungen} abgelegt.

Das Skript basiert auf vielen verschiedenen Quellen, welche falls möglich referenziert werden.

Das Skript wird von der Lehrperson geschrieben und auf dem aktuellsten Stand gehalten.
Selbstverständlich sind Fehler vorhanden, falls Sie einen bemerken, wird ein Hinweis darauf
dankend entgegengenommen.